
\begin{tikzpicture}[x=12mm,y=12mm,font=\small]
\draw[otMuted,->](-2.5,0)--(2.6,0)node[right]{$x$};\draw[otMuted,->](0,-2.4)--(0,2.5)node[above]{$y$};
\draw[otNavy](0,0)circle(2);\draw[otMuted,dashed](1,-2.2)--(1,2.2);
\fill[otNavy](-2,0)circle(1.2pt)node[below left]{$A$};\fill[otNavy](2,0)circle(1.2pt)node[below right]{$B$};
\fill[otNavy](1,1.7320508)circle(1.2pt)node[above right]{$M_+$};
\fill[otNavy](1,-1.7320508)circle(1.2pt)node[below right]{$M_-$};
\geoSolution{\draw[otGreenDark](-2,0)--(1,1.7320508)--(2,0);
\draw[otGreenDark](-2,0)--(1,-1.7320508)--(2,0);
\draw[otGreenDark](.8700962,1.6570508)--(.9450962,1.5271470)--(1.075,1.6021470);
\draw[otGreenDark](.8700962,-1.6570508)--(.9450962,-1.5271470)--(1.075,-1.6021470);}
\node[below left]at(0,0){$O$};
\end{tikzpicture}
\par\smallskip\legende{Le cercle de diamètre $[AB]$ et une droite verticale ; l’exemple représenté correspond à $k=20$.}
